\chapter{Tata Kelola AI}
Penggunaan AI yang serabutan, asal-asalan, mulai menimbulkan masalah.
Ada kebocoran data perusahaan karena karyawan memasukkan data nasabah
ke mesin AI (chat, LLM) tanpa sadar bahwa data tersebut digunakan juga
oleh penyedia jasa AI untuk menjadi bagian dari data pelatihan.
Untuk mengatasi hal ini maka harus mulai diperhatikan aspek tata kelola
(governance) dari AI.

Tata kelola secara umum menyangkut lima (5) pertanyaan berikut.
\begin{itemize}
\item Siapa (who)
   yang diperkenankan mengembangkan, melakukan penyetujuan,
   dan menggunakan AI;
\item Sistem AI mana 
   (what) yang boleh digunakan dan untuk keperluan apa;
\item Data yang mana (which)
   yang dapat diberikan kepada sistem AI;
\item Bagaimana (how) kita dapat memastikan bahwa AI aman,
   fair, dan handal (reliable);
\item Siapa (who) yang bertanggungjawab (accountable)
   jika sistem AI melakukan kesalahan dan bahkan
   mencelakakan manusia.
\end{itemize}

\section{Siapa Yang Boleh Menggunakan AI}
Siapa yang boleh menggunakan AI? Jawabannya hampir pasti semuanya.
Namun maksudnya bukan seperti itu.
Di dalam konteks perusahaan, siapa yang boleh menggunakan AI untuk
mengelola atau memproses data perusahaan.
Ada banyak data yang bersifat rahasia (\emph{confidential}) yang tidak
boleh diberikan ke pihak lain. Tidak boleh diberikan ke mesin AI.
Pertanyaan tentang siapa yang boleh itu dikaitkan dengan hal ini.

Masalahnya adalah jika perusahaan atau instansi tidak memberikan akses
kepada mesin / layanan AI, makan karyawan akan membuat atau menggunakan
mesin AI sendiri. Ini yang disebut \emph{shadow AI}.
AI yang tidak resmi digunakan tetapi ada dalam perusahaan.
Keberadaan shadow AI inilah yang berisiko.
